%---------------------------------------------------
% کلاس سند: article برای اسناد متنی، با اندازه کاغذ A4 و فونت پایه 16
\documentclass[a4paper,16pt]{article}

% بسته‌های ریاضی برای معادلات و نمادهای پیشرفته
\usepackage{amsmath, amssymb, amstext}

% بسته رنگ برای استفاده از رنگ در متن یا محیط‌ها
\usepackage{xcolor}

% بسته گرافیک برای درج تصاویر
\usepackage{graphicx}

% تنظیمات حاشیه صفحه
\usepackage{geometry}

% بسته لینک‌دهی (هایپرلینک‌ها)
\usepackage{hyperref}

% بسته صفحه‌آرایی برای تنظیم هدر و فوتر
\usepackage{fancyhdr}


% بسته xepersian برای پشتیبانی از زبان فارسی با XeLaTeX


% تعیین فونت فارسی (فایل فونت باید موجود باشد یا در سیستم نصب باشد)
\usepackage{xepersian}
\settextfont{B Nazanin}

\setlatintextfont{Times New Roman}

% تنظیم اندازه حاشیه‌ها از هر طرف برابر با 1 اینچ
\geometry{margin=1in}

% تنظیم سبک fancy برای سربرگ و پاورقی
\pagestyle{fancy}

% سمت چپ هدر: مشخص‌کننده سری تمرین و نام درس
\lhead{تمرین سری سوم درس مکانیک تحلیل }

% وسط هدر: نام دانشکده
\chead{دانشکده فیزیک دانشگاه صنعتی خواجه نصیرالدین طوسی}

% سمت راست هدر: تاریخ روز به‌صورت خودکار
\rhead{\today}

% پایین صفحه: شماره صفحه
\cfoot{\thepage}

% شروع سند
\begin{document}
	
	% صفحه عنوان (بدون هدر و فوتر)
	\thispagestyle{empty}
	\begin{center}
		% لوگوی دانشگاه (باید فایل تصویر در مسیر پروژه باشد)
		\includegraphics[width=0.5\textwidth]{../../../images/image-E_3/image-1.png} \\[10pt] 
		
		% عنوان سند
		\textbf{\LARGE عنوان:تمرین سری‌سوم}\\[20pt]
		
		% نیم‌سال تحصیلی
		\textbf{\LARGE نیم‌سال تحصیلی:4041 }\\[10pt]
		
		% نام مدرس درس
		\textbf{\Large مدرس: دکتر محمد انصاری‌فرد}\\[10pt]
		
		% مبحث تمرین 
		\textbf{\Large مبحث تمرین:مکانیک نیوتونی   }\\[10pt]
		
		% مهلت تحویل تمرین
		\textbf{\Large مهلت تحویل:۱۲آبان }
	\end{center}
	
	% صفحه جدید برای ادامه سند
	\newpage
	
	% ایجاد فهرست مطالب به‌صورت خودکار (براساس \section و \subsection و ...)
	\tableofcontents
	\newpage
	
	
	\section{سوال اول}

ذره‌ای به جرم \( m \) از سطح شیبداری تحت تأثیر نیروی گرانش به پایین می‌لغزد.  
اگر نیروی مقاومی در برابر حرکت آن وجود داشته باشد که برابر است با
\[
f = k m v^2
\]
که در آن \( k \) عددی ثابت است،  
نشان دهید که زمان لازم برای این‌که ذره از حالت سکون مسافت \( d \) را طی کند، برابر است با:
\[
t = \frac{1}{\sqrt{k g \sin\theta}} \cosh^{-1}\!\left(e^{k d}\right)
\]
که در آن \( \theta \) زاویه شیب سطح شیبدار است.

	\section{سوال دوم}
	
	سرعت ذره‌ای به جرم \( m \) به‌صورت زیر داده شده است:
	\[
	v = \frac{\alpha}{x}
	\]
	که در آن \( x \) تغییر مکان ذره است.  
	نیروی عامل این حرکت را بر حسب \( x \) و مشتقات مناسب آن بیابید.
	
	
	
	\section{سوال سوم}

اگر پرتابه‌ای با سرعت اولیه $v_0$ و در جهتی که با افق زاویه $\alpha$ می‌سازد از مبدأ دستگاه مختصات شلیک شود، زمان لازم را بیابید تا پرتابه خطی را قطع کند که از مبدأ می‌گذرد و با افق زاویه $\beta$ می‌سازد.

	\section{سوال چهارم}
	
	دو قطعه با جرم‌های نامساوی \(m_1\) و \(m_2\) به وسیله ریسمانی که از روی یک قرقره صاف و بدون اصطکاک می‌گذرد به هم وصل شده‌اند هر قطعه روی یک سطح شیبدار قرار دارد که زاویه شیب هر دو برابر \(\theta\) فرض می‌شود. ضریب اصطکاک جنبشی بین هر جسم و سطح مربوطه برابر \(\mu_k\) است. 
	\\ 
	مشخص کنید برای حرکت با سرعت ثابت زاویه \(\theta\) باید چه مقداری داشته باشد.  
	


	
	\section{سوال پنجم}
	یک موشک در فضای خلأ (بدون مقاومت هوا) ابتدا ساکن است و سوخت را با نرخ ثابتِ \(\dot{m}\) می‌سوزاند و گازها را با سرعتِ ثابت \(u\) نسبت به موشک به بیرون می‌پرانَد. جرم اولیه موشک (شامل سوخت) را \(M_0\) و جرم نهایی پس از سوزاندنِ تمام سوخت را \(M_f\) در نظر بگیرید.
	
	\begin{enumerate}
		\item نشان دهید سرعت موشک در انتهای مرحله رانش برابر است با
		\[
		v = u \ln\!\left(\frac{M_0}{M_f}\right)
		\]
		\item به‌طور فیزیکی توضیح دهید چرا با کاهش جرمِ موشک، شتاب آن افزایش می‌یابد.
	\end{enumerate}
	
	
	\vspace{10pt}
	\textbf{موفق باشید.}
\end{document}
